\documentclass{article}
\usepackage{graphicx}
\usepackage{float}
\setlength{\parindent}{0pt}
\usepackage{caption}
\usepackage{fancyvrb}
\usepackage{tcolorbox}
\usepackage[a4paper, total={6in, 8.5in}]{geometry}
\setlength{\footskip}{80pt}
\usepackage{tabularx}
\usepackage{booktabs}
\usepackage{enumitem}

\title{Project - Protocol Hardening}
\author{Aleksander Kłap}
\date{08.06.2026}

\begin{document}

\maketitle
\pagebreak

\section{Protocol Hardening}
The redesign (Chapter 7) controls where each protocol may flow. 
This chapter defines how each protocol is secured on the device and conduit level.

\subsection{Profinet}
Profinet has no authentication or encryption features, so it cannot be hardened at the 
protocol level. Instead it is secured by isolation.
\begin{itemize}
    \item Layer 2 only - Profinet stays inside each OT Production station (VLAN 101 - 106).
        It is never routed, and never crosses through firewall. Even if routing would be 
        executed (through the attack) router would drop the communication.
    \item No inter-station Profinet communication - Each PLC talks Profinet (with I/O, HMI) only 
        within own VLAN.
    \item DCP Set disabled - desible DCP Set through TIA Portal for devices that support this feature.
        This prevents rename/reconfigure attacks through DCP.
\end{itemize}

\subsection{OPC UA}
OPC UA is the only routed OT data protocol. Hardening is specified in SR-03. \textit{(see Sercurity Requirements 6.1 TODO)}
\begin{itemize}
    \item Security mode Sign\&Encrypt - Security policy Basic256Sha256. Mode 'none' is disabled
        on all PLC and Festo Controllers servers.
    \item Certificate authentication - Clients (OPC UA proxy, Cobot/MLvisionPC) authenticate with application
        certificates. Annonymous and username/password sessions are disabled.
    \item Scoped access - The Cobot/MLvisionPC can reach only PLC2 on Station 2, and no other PLC's.
\end{itemize}

\subsection{MQTT}
MQTT is used by the Festo CXV device for telemetry. It's hardening is specified in SR-04 \textit{(see Sercurity Requirements 6.1 TODO)}
\begin{itemize}
    \item MQTT broker in IDMZ - The MQTT CXV device publishes only to a dedicated MQTT broker placed in IDMZ. IT nodes 
        can only subscribe to the broker. No direct OT-IT MQTT path exists.
    \item TLS 1.2+ - TLS encryption enforced on all MQTT connections.
    \item Client certificates - MQTT CXV device must authenticate with certificate for publish connection to the broker. No annonymous connections allowed.
    \item Topic ACL's - THe MQTT CXV device may publish only to it's own topics. IT nodes (MES) cannot publish to the OT topics.
\end{itemize}

\subsection{HTTP/Web Interfaces}
Some OT devices need to expose plaintext HTTP web interfaces. Hardening is specified in SR-05 \textit{(see Security Requirements 6.1 TODO)}

\begin{itemize}
    \item Disabled by default - All HTTP servers are disabled unless strictly needed for operational functionality.
    \item Reverse proxy - Where web access is needed (e.g. MES reading the PLC alarm table), the MES connects to the Nginx reverse proxy in the IDMZ over HTTPS/TLS 1.2+. Nginx terminates the connection, fetches the OT-side resource, and returns it. Plaintext HTTP stays inside the OT zone.
    \item Access restriction - Ngnix authenticate requests, logs them and blocks any OT server that should not be reached from IT zone.
        Firewall rules restrict web access to the Nginx TCP/IP port only.
\end{itemize}

\subsection{S7comm+ Engineering Access}
S7comm+ connection is used only for programming PLC's and HMI's. Hardening is sepcified in SR-06 and SR-07 TODO link to SR's
\begin{itemize}
    \item Single source - S7comm+/TCP102 to any PLC or HMI is allowed only from the dedicated Engineering Workstation in VLAN40.
    \item Conduit controll - Engineering Worksation S7comm+ connection initialization passes through the firewall, enforcing restriction of Engineering Station to PLC's and HMI's only.
    \item User Authentication - PLC's and HMI's access requires individual login/password authentication with role separation (implemented in TIA portal).
\end{itemize}

\end{document}

